\subsection{IND-CPA Public-Key Encryption}
\label{app:ind-cpa}

We use the single-challenge experiment of Bellare et al.\
\cite[full version, Definition~2.1]{BellareEtAl1998EncryptionNotions}
with the encryption syntax of Section~\ref{sec:ind-cpa}. Correctness is
as in Appendix~\ref{app:nce}, with deterministic decryption.
For $b\in\bits$, the experiment $\CPAGame^b_{\mathcal A}(1^\lambda)$
samples $(\pk,\sk)\sample\KeyGen(1^\lambda)$, gives
$(1^\lambda,\ell,\pk)$ to $\mathcal A$, and offers these commands:
\begin{description}
  \item[$\Challenge(m_0,m_1)$.]
    For $m_0,m_1\in\bits^\ell$, return
    $c\sample\Enc(\pk,m_b)$ using fresh randomness.
    At most one challenge is accepted.
  \item[$\Finish(b')$.]
    Terminate and output $b'\in\bits$.
\end{description}
Malformed commands and further challenges return $\bot$ without changing
state.  Termination without an output bit is interpreted as $\Finish(0)$.
The messages may depend on the public key.  No encryption oracle is
needed: the adversary can run $\Enc(\pk,\cdot)$ itself, both before
and after the challenge.

\begin{definition}[IND-CPA security]
\label{def:ind-cpa}
The encryption scheme is IND-CPA secure if, for every efficient adversary
$\mathcal A$, there is a negligible function $\varepsilon_{\mathcal A}$
such that
\[
  \left|
    \Pr[\CPAGame^0_{\mathcal A}(1^\lambda)=1]
    -\Pr[\CPAGame^1_{\mathcal A}(1^\lambda)=1]
  \right|\leq\varepsilon_{\mathcal A}(\lambda).
\]
\end{definition}

The negligible security bound is uniform over polynomially bounded
public message lengths, as in Definition~\ref{def:nce}.
